% Pursuing Stacks (cotes 134-2 to 134-8): archive cote 134-7, site page 40 (PDF page 41), typescript page 482.
% typescript. Draft compared: ps6.tex:1935-2003.
% AI-transcribed by Claude Opus 5.5 (claude-opus-5-5, adaptive thinking, effort high), Message Batches API, batch ps3 (msgbatch_01XdFmuVVBAcgc7EFxdtCqsj), 2026-11-12; not yet checked by eye.
\scanpage{134-7/40}
\mspage{482}

associating to any affine scheme \inserted{S=} Spec(k') over k the ring of sections
of  it\deleted{'}s usual Zariski structure sheaf,  which ring is canonically isomor-%
phic to k' itself !  The operation (1) is the usual linearization opera-%
tion with respect to this sheaf of rings,  working in the topos of
fpqc sheaves of sets over k which we described at some length in section
111  (p.447).  As I was fearing that working in such a thing would cause
anguish to a number of prospective readers,  I took pains to translate
unipotent bundles from the geometric language which is the suggestive
one,  to the language of commutative algebra which is more liable to hide
than to disclose geometrical meaning;  so much so that in the process I
m\xout{x}yself lost contact somewhat with the geometric flavor,  and more speci-%
fically still with this basic fact,  that in our context of unipotent
bundles and complexes of such,  the ``natural'' coefficients for cohomology
\margin{$X(n)_*,\hand{\pi}_{n+1}$} (such as the $H^{n+2}(\unsure{X(n)_*},\hand{\pi}_{n+1})$ groups occuring in Posthikov dévissage) are
by no means``\uline{discrete}''\hand{/}one's such as $\underline{Z}$ or k ,  but \uline{quasi}-\uline{coherent} one's
namely provided by quasi-coherent sheaves of \xout{modules} $\underline{O}_k$-modules or com-%
plexes of such.  Thus in the above Postnikov obstruction group, $\hand{\pi}_{n+1}$
does \uline{not} stand as a constant group of coefficients (if it was,  this
would drag us into the niceties and difficulties of etale cohomology
for the components $X(n)_i$ of the semisimplicial unipotent bundle $X(n)_*$);
but using the k-module structure of $\hand{\pi}_{n+1}$ for defining a \uline{quasi}-\uline{coherent}
\unsure{\hand{s}}heaf of modules  $W(\hand{\pi}_{n+1})$  ``over k'' i.e.  over $Aff_{/k}$ ,  it is this
``continuous'' sheaf (or ``vector bundle'') over k, lifted of course to the
various components \xout{\unsure{at}} $X(n)_i$ ,  which yields the correct answer.  This
was kind of clear in my mind the very first day when I started reflecting
on schematization, \inserted{even} before \deleted{even} introducing formally unipotent bundles
(pages 443-444),  but this instinctive understanding later became dulled
somewhat,  largely due it seems to me to the concession I had made to
\xout{commutative} algebra,  giving up to some extent the language of geometry.

Let's recall that the operation (1) may be defined as the solution
of a universal problem,  namely sending the non-linear object X into
a ``linear'' one namely into a sheaf of $\underline{O}_k$-modules (or a $\underline{O}_k$-Module,  as
we'll simply say).  This is expressed by a canonical \xout{\unsure{map}} map of sheaves
of sets
\[ (3) \qquad X \longrightarrow \underline{O}_k^{(X)} \]
(which I am tempted to call the ``\uline{exponential}'' \uline{map} for X, and denote by
a corresponding symbol such as $exp_X$ ) ,  giving rise,  for every Module
(over $\underline{O}_k$)  to a corresponding map which is \uline{bijective}

\margin{40}
